\par\medskip
\noindent
\underline{1. Modelos de Néron de los esquemas abelianos: notación; el
emparejamiento canónico}
\[
 \Phi_o\times\Phi'_o\longrightarrow(\mathbb Q/\mathbb Z)_k,
\]
\underline{y la biextensión canónica $W^o$ de $(A^o,{A'}^o)$.}

\medskip
\noindent
1.0. A lo largo de esta exposición fijamos un rasgo $S$, para el cual
utilizamos las notaciones habituales $s$, $\eta$, $k$, $K$, etcétera,
de I~1. Sea $A_K$ un esquema abeliano sobre $K$, y denotemos por $A'_K$
su esquema abeliano dual [5]. La dualidad entre $A_K$ y $A'_K$ viene
dada por una correspondencia divisorial $W_K$ sobre $(A_K,A'_K)$, es
decir, por un haz invertible sobre $A_K\times A'_K$, birigidificado con
respecto a las secciones identidad de $A_K$ y $A'_K$. Queda determinado,
salvo isomorfismo único, por su clase (cf. VIII~3.1)
\begin{equation}\tag{1.0.1}
 w_K\in\mathrm{Biext}^1(A_K,A'_K;\mathbb G_{mK}).
\end{equation}
Desempeñará un papel importante el emparejamiento correspondiente
(VIII~2.2)
\begin{equation}\tag{1.0.2}
 \varphi:T_\ell(A_K)\times T_\ell(A'_K)
 \longrightarrow\mathbb Z_\ell(1)_K
 \overset{\mathrm{def}}{=}T_\ell(\mathbb G_{mK}),
\end{equation}
donde $\ell$ es un número primo arbitrario. Cuando
$\ell\ne p=\operatorname{char}k$, dar $\varphi$ equivale a dar un
emparejamiento de $\mathbb Z_\ell$-módulos libres de tipo finito
\begin{equation}\tag{1.0.3}
 \varphi:T_\ell(A_K(\overline K))\times
 T_\ell(A'_K(\overline K))\longrightarrow T_\ell(\overline K^*),
\end{equation}
compatible con la acción de
\[
 \pi=\operatorname{Gal}(\overline K/K)
\]
sobre los tres módulos en cuestión, donde $\overline K$ es una clausura
separable de $K$. Se sabe (VIII~3.2) que este emparejamiento es una
dualidad perfecta: identifica cada uno de los módulos
$T_\ell(A_K(\overline K))$ y $T_\ell(A'_K(\overline K))$ con el dual
del otro, con valores en el módulo invertible
$\mathbb Z_\ell(1)=T_\ell(\overline K^*)$.

En ocasiones supondremos elegida una polarización de $A_K$; así, se
nos da un elemento
\begin{equation}\tag{1.0.4}
 \xi\in\operatorname{NS}(A_K/K)
 =\underline{\operatorname{NS}}_{A_K/K}(K),
\end{equation}
donde
\[
 \underline{\operatorname{NS}}_{A_K/K}
 \overset{\mathrm{def}}{=}
 \underline{\operatorname{Pic}}_{A_K/K}/
 \bigl(\underline{\operatorname{Pic}}{}^o_{A_K/K}=A'_K\bigr).
\]
De la manera habitual, este define un homomorfismo
\begin{equation}\tag{1.0.5}
 \psi_\xi:A_K\longrightarrow A'_K,
\end{equation}
por tanto un homomorfismo
\begin{equation}\tag{1.0.6}
 T_\ell(\psi_\xi),\quad\text{o simplemente}\quad
 \psi_\xi:T_\ell(A_K)\longrightarrow T_\ell(A'_K),
\end{equation}
y una forma bilineal $\varphi_\xi$
\begin{equation}\tag{1.0.7}
 \begin{split}
 \varphi_\xi(x,y)&=\varphi(x,\psi_\xi(y)),\\
 \varphi_\xi:T_\ell(A_K)\times T_\ell(A_K)
 &\longrightarrow\mathbb Z_\ell(1)_K,
 \end{split}
\end{equation}
que, para $\ell\ne p$, puede interpretarse también como un
emparejamiento de $\mathbb Z_\ell$-módulos libres de tipo finito con
operadores:
\begin{equation}\tag{1.0.8}
 \varphi_\xi:T_\ell(A_K(\overline K))\times
 T_\ell(A_K(\overline K))\longrightarrow T_\ell(\overline K^*).
\end{equation}
Estas construcciones tienen sentido para todo elemento del grupo de
Néron--Severi de $A_K$, y producen una forma \underline{alternante}
$\varphi_\xi$, en cualquiera de las interpretaciones (1.0.7) o (1.0.8)
(VIII~2.2.11). El hecho de que $\xi$ sea una \underline{polarización}
implica que $\psi_\xi$ en (1.0.5) es una isogenia; equivalentemente,
puesto que $\dim A_K=\dim A'_K$, es un epimorfismo de grupos. Esto
equivale también a la inyectividad de $T_\ell(\psi_\xi)$ en (1.0.6)
(o a que tenga conúcleo finito) y, finalmente, a que la forma
$\varphi_\xi$ de (1.0.7) sea \underline{separante}.

\medskip
\noindent
1.1. Recordemos que existe un \underline{esquema liso} $A$ sobre $S$,
llamado \underline{modelo de Néron de $A_K$}, que representa el funtor
\begin{equation}\tag{1.1.1}
 S'\longmapsto\operatorname{Hom}(S'_\eta,A_K)
\end{equation}
sobre la categoría de los $S$-esquemas lisos. Así, por definición,
existe un isomorfismo, funtorial en el $S$-esquema relativo liso $S'$,
\begin{equation}\tag{1.1.2}
 \operatorname{Hom}_S(S',A)\simeq
 \operatorname{Hom}_K(S'_\eta,A_K).
\end{equation}
Néron [19] construyó $A$ cuando el cuerpo residual $k$ de $S$ es
perfecto. Raynaud extendió la construcción al caso general en un
seminario no publicado celebrado en el IHES en 1966--67 (véase [21]
para el enunciado de los principales resultados de Raynaud, que
completan los de Néron en loc. cit.). Admitiremos aquí este resultado
de existencia, junto con el hecho de que $A$ es un \underline{esquema}
cuasiproyectivo sobre $S$. Aplicando (1.1.2) cuando $S'$ es liso
\underline{sobre} $\eta$, se obtiene también que el siguiente
morfismo canónico es un isomorfismo:
\begin{equation}\tag{1.1.3}
 A\times_S\operatorname{Spec}(K)\xrightarrow{\ \sim\ }A_K.
\end{equation}
Esto justifica la notación emparejada $A,A_K$; cuando convenga,
escribiremos también $A_\eta$ en lugar de $A_K$. Utilizando (1.1.3)
como isomorfismo de identificación, la biyección canónica (1.1.2) es
simplemente la aplicación dada por la \underline{restricción a las
fibras genéricas}.

Puesto que el funtor (1.1.1) posee manifiestamente una estructura de
grupo conmutativo, el modelo de Néron $A$ que lo representa posee a su
vez una estructura de \underline{esquema en grupos conmutativo liso
sobre $S$}. El resto de esta exposición está dedicado al estudio de
este esquema en grupos.

Un invariante importante procedente del modelo de Néron es, por
supuesto, el esquema en grupos conmutativo liso de tipo finito sobre
$k$
\begin{equation}\tag{1.1.4}
 A_o=A\times_Ss=A_k.
\end{equation}
Conforme a la notación general, denotamos por $A_o^o$ su componente
neutra y consideramos también el esquema en grupos
\begin{equation}\tag{1.1.5}
 \Phi_o=A_o/A_o^o
\end{equation}
de componentes conexas de $A_o$, que es finito y étale sobre $k$. En
términos de una clausura separable $\overline k$ de $k$, queda por tanto
determinado por el grupo finito ordinario
\begin{equation}\tag{1.1.6}
 \Phi_o(\overline k),
\end{equation}
junto con la acción natural de
\[
 \pi_o=\operatorname{Gal}(\overline k/k)
\]
sobre este grupo.

Junto con el modelo de Néron $A$ de $A_K$, estudiaremos
\begin{equation}\tag{1.1.7}
 A'=\text{el modelo de Néron de }A'_K,
\end{equation}
que es asimismo un esquema en grupos conmutativo liso de tipo finito
sobre $S$. Denotaremos por
\begin{equation}\tag{1.1.8}
 \Phi'_o=A'_o/{A'}_o^o
\end{equation}
el grupo de componentes conexas de su fibra especial
\sisi{$A_o$}{$A'_o$}.
